\ficha{Bordo e Schwartz (1994), Specie Standard as a Contingent Rule}{bordoschwartz1994}

\begin{identificacao}
BORDO, Michael D.; SCHWARTZ, Anna J. \textit{The specie standard as a
contingent rule}: some evidence for core and peripheral countries, 1880-1990.
Cambridge, MA: National Bureau of Economic Research, 1994. 69 p. (NBER Working
Paper, n. 4860). DOI: 10.3386/w4860.

Working paper, lido integralmente (44 p. de texto, mais notas, Tabela 1 e
apêndice de dados), em inglês. A versão lida é a do NBER, setembro de 1994.
\end{identificacao}

\bloco{Problema e tese}
O artigo pergunta por que alguns países sustentaram a conversibilidade metálica
entre 1880 e 1914 e outros não, e o que a adesão significava para quem
emprestava. A resposta é que o padrão metálico anterior a 1914 não era um
automatismo, e sim uma regra contingente: um compromisso de manter fixo o preço
da moeda em espécie, com cláusula de escape para emergência exógena bem
compreendida, tipicamente a guerra, e com a obrigação de retomar depois na
paridade original. Cumprir a regra funcionava como sinal de compromisso crível
e dava ao governo acesso a senhoriagem e a crédito em condições melhores. Os
países do centro cumpriram a regra; boa parte da periferia, e a América Latina
em particular, não.

\bloco{Argumento}
\begin{enumerate}[leftmargin=*,nosep]
\item A regra é tratada como mecanismo de compromisso contra inconsistência
temporal, e não como automatismo impessoal da tradição de Simons. Sem
compromisso crível, o governo tem incentivo a surpreender com inflação e a dar
calote implícito na dívida, e os credores antecipam isso (pp. 3-4).
\item O conteúdo operacional da regra é estreito. Suspender em guerra e voltar
depois na paridade antiga é adesão. Discrição é adiar a retomada além do prazo
razoável, ou fixar de novo em paridade desvalorizada (p. 4). A regra precisa
ser transparente e ter poucas contingências, porque crise financeira e choque
de termos de troca são causas mais difíceis de verificar que guerra, e invocar
a cláusula de escape nesses casos gera suspeita de que se está exercendo
discrição (p. 5).
\item O cumprimento é sustentado por reputação e, em alguns casos, por
dispositivo legal, como a garantia constitucional sueca de conversibilidade
(p. 6). No plano internacional, os autores listam quatro mecanismos: acesso ao
mercado de capitais do centro, as regras do jogo, a hegemonia inglesa e a
cooperação episódica entre bancos centrais (p. 7). O sistema é descrito como
assimétrico: a Inglaterra usava a taxa bancária para manter a conversibilidade
e os demais aceitavam o ditame das paridades fixas, deixando a oferta de moeda
responder passivamente (p. 8).
\item A Tabela 1 organiza a cronologia de conversibilidade e suspensão de 21
países em sete grupos. Sob o padrão-ouro clássico o registro do centro é
descrito como impecável (p. 14), e Austrália, Canadá, escandinavos, Europa
ocidental menor e Japão aparecem tão fiéis quanto o centro (p. 31). Argentina, Brasil e
Chile não, e os elementos comuns apontados são guerra, ameaça de guerra e
política fiscal e monetária incompatível com câmbio fixo (p. 18).
\item A evidência sobre fluxos de capital compara chamadas de capital sobre
novas emissões em Londres com a cronologia de adesão de Argentina e Estados
Unidos, 1865-1914. Os autores tratam o resultado como sugestivo, não como teste
(p. 33). O caso argentino fecha em 1899 com a restauração da conversibilidade e
a criação da Caja de Conversión, que amarrou as mãos da autoridade monetária,
mas o texto reconhece que o determinante principal do fluxo foi a abertura
econômica e a estabilidade política, não a regra (p. 34).
\item A seção estatística compara quatro regimes entre 1880 e 1989 em inflação,
crescimento monetário e déficit público sobre PIB, e estima choques de oferta e
demanda por VAR bivariado (p. 41). Em todos os regimes o déficit médio é maior na
periferia que no centro (p. 39), e nos países latino-americanos os choques de
demanda superam os de oferta em quase todos os regimes, o que os autores leem
como uso de política discricionária (p. 43).
\item A conclusão explica a diferença entre centro e periferia por estágio de
desenvolvimento: o centro já estava industrializado, tinha renda per capita
crescente, sistemas políticos estáveis e grupos dominantes que percebiam o benefício da
estabilidade monetária, e como economia diversificada estava menos exposto a
oscilações de preço de produto primário (p. 44).
\end{enumerate}

\bloco{Conceitos e definições operacionais}
Regra contingente: compromisso de preço fixo da moeda em espécie com cláusula
de escape acionável em emergência exógena bem compreendida, mais obrigação de
retomar na paridade original (p. 4). Discrição: adiamento da retomada ou
fixação em paridade desvalorizada (p. 4). Credibilidade: propriedade observável
pelo comportamento de câmbio e juros dentro dos pontos de ouro, e não
declarada; Giovannini é citado para o período 1899-1909 (p. 6). Centro e
periferia: classificação por grupo, com França, Alemanha, Reino Unido e Estados
Unidos no centro; a inclusão dos Estados Unidos é discutida e defendida
(pp. 12 e 14). Currency board aparece de forma operacional, não teórica, como
arranjo que ``em essência amarra as mãos da autoridade monetária'', descrito
para a Caja de Conversión (p. 34). Selo de boa conduta: a adesão como
sinalização de probidade financeira aos credores de Londres (p. 31).

\chave{In 1906 Brazil restored convertibility to prevent continued appreciation
of the milreis exchange rate that was harmful to coffee and rubber
exporters}{p. 17}

\bloco{Evidência e método}
Três camadas. Primeira, uma cronologia narrativa de 21 países montada sobre
literatura secundária nacional, com Fritsch e Franco (1992) e Peláez e Suzigan
(1976) para o Brasil, Cortés Conde para a Argentina, Llona Rodríguez para o
Chile, entre outros; não há fonte primária nacional. Segunda, séries anuais de
chamadas de capital sobre novas emissões em Londres, cedidas por Lance Davis e
originadas do estudo de Davis e Huttenback sobre o imperialismo britânico,
apresentadas em gráfico com as fases de suspensão sombreadas (p. 32).
Terceira, estatística descritiva de 1880 a 1989 para inflação medida por
deflator do PIB, crescimento monetário e razão déficit sobre PIB, mais VAR
bivariado em preço e produto com as restrições de Blanchard e Quah, dois lags e
três subperíodos (p. 41). Não há contrafactual formal nem teste de
significância da adesão.

\bloco{Agentes e vertentes}
Bordo e Schwartz escrevem contra duas posições: a crítica de De Cecco, para
quem o padrão-ouro foi favorável ao centro e desfavorável à periferia, e a
crítica keynesiana de subordinação da estabilidade interna à conversibilidade
externa (p. 1). Dialogam com Eichengreen, Giovannini, Fishlow, Frieden,
Gallarotti e Bloomfield. A nota 1 registra o debate entre escola da moeda e
escola bancária como antecedente da ideia de regra, com a escola da moeda
defendendo emissão fiduciária variando com a reserva de ouro. O Brasil aparece
sem nenhum agente nomeado: não há Campista, Bulhões, Murtinho nem Rui Barbosa,
e o Convênio de Taubaté não é mencionado.

\bloco{Críticas e limites}
A evidência sobre capital é frágil pelo próprio reconhecimento dos autores, que
pedem pesquisa mais sistemática (p. 35). Duas passagens andam contra a tese: o
forte retorno dos fluxos para a Argentina entre 1886 e 1889 sob moeda
inconversível, que os autores atribuem à aposta na prosperidade de longo prazo
(p. 34), e a nota 14, que registra que o Chile tomou empréstimo em libras a
taxas correntes durante a inconversibilidade. A identificação de choques por
VAR é apresentada com ressalva explícita de limitação. Há circularidade
admitida: o texto não separa se o desempenho macroeconômico explica a adesão ou
se a adesão explica o desempenho (p. 40).

Para o caso brasileiro há um limite mais específico. A régua do próprio artigo
classifica como discrição a fixação em paridade desvalorizada, e a Caixa de
Conversão fixou 15 pence, muito abaixo do par legal de 27 pence, mas a Tabela 1
registra 1906 a 1914 como período de conversibilidade. A tipologia binária
adesão contra não adesão não distingue retorno à paridade antiga de
estabilização em taxa nova, que é justamente o eixo do debate brasileiro.

\chave{for the Latin American peripheral countries, not only the aftermath of
war but deflation generally were reasons for absence of commitment to
convertibility}{p. 19}

\begin{dialogo}
\item Procurar nos jornais o argumento de sinalização a credores: crédito no
exterior, confiança dos capitalistas, praça de Londres, condições do próximo
empréstimo. Testa se o selo de boa conduta (p. 31) era categoria nativa dos
agentes de 1906 ou reconstrução posterior da literatura.
\item Procurar o argumento do par legal, 27 pence, contra a taxa de
estabilização, 12 ou 15 pence. Testa se a lógica de retomada na paridade
original, que Bordo e Schwartz tomam como definidora da regra (p. 4), tinha
defensores no debate brasileiro e em que momento perdeu.
\item Procurar como os jornais qualificam a suspensão do troco em 1914. Se a
tratam como medida legítima imposta pela guerra e provisória, é indício de que
o vocabulário da cláusula de escape circulava; se a tratam como falência do
arranjo, o vocabulário não circulava.
\item Procurar menção à Caja de Conversión argentina e ao caso argentino de
1899. Testa a difusão do argumento de amarrar as mãos da autoridade monetária
(p. 34) como modelo institucional invocado no Brasil.
\item Procurar a defesa da Caixa formulada como proteção a café e borracha
contra apreciação. Testa se a leitura da p. 17, que faz do retorno de 1906 uma
defesa cambial de exportador e não uma busca de disciplina, corresponde ao que
os agentes diziam publicamente.
\end{dialogo}

\usonocapitulo{Parte I. Serve para três afirmações. Primeira, fornece a
definição operacional de regra contingente e de cláusula de escape que
organiza a discussão de credibilidade no capítulo, com o critério de adesão e
discrição enunciado na p. 4. Segunda, dá a formulação canônica da assimetria
centro-periferia e sua explicação por estágio de desenvolvimento (p. 44), que
o capítulo confronta com o caso brasileiro. Terceira, registra o modo como a
literatura internacional lê a Caixa de Conversão sem consultar o debate
brasileiro, tratando o retorno de 1906 como defesa cambial de exportadores de
café e borracha (p. 17), o que abre espaço para a contribuição da dissertação.
Dois trechos servem a outras partes: a exigência de política contracionista
como condição do funding loan de Londres (p. 17) entra na Parte II, e a
descrição da Caja de Conversión como quase currency board (p. 34) entra na
Parte IV.}

\chave{pegging to specie at a devalued parity}{p. 4}
